% ======================================================================
%  第 4 讲  紧化：四维的桥塌成二维的镜像
%  页锚以 B 版原文（\cite{kw06}）的印刷页为准
% ======================================================================
\section{紧化：四维的桥塌成二维的镜像}
\label{sec:compact}

\begin{goals}
\begin{enumerate}
\item 说清楚把理论放在 $\Sigma\times C$ 上、让 $C$ 变小时，低能理论为什么是以 $\MH$ 为靶空间的二维 sigma 模型。
\item 说清楚这个二维理论在参数特殊值处是 A-模型还是 B-模型，以及 $\Psi$ 如何统一描述它们。
\item 说清楚四维的 $S$-对偶为什么变成二维的镜像对称，并把核心交换图的两条路径对上。
\end{enumerate}
\end{goals}

\subsection{把世界做成乘积 $\Sigma\times C$}

本讲对应 \cite{kw06} 的 §4--§5。几何 Langlands 的语言是二维的：曲线 $C$ 上的层、Hecke 算子、平坦联络、$\Dmod$-模。四维的 Kapustin--Witten 理论要「长出」一个二维世界，第一步就是紧化。原文的选择是把四维理论放在
\begin{equation}
M=\Sigma\times C
\label{eq:M-prod}
\end{equation}
上，其中 $C$ 是做几何 Langlands 的那条紧 Riemann 面（亏格 $g>1$，保证一般平坦联络不可约），$\Sigma$ 是另一个二维流形，充当低能二维理论的「时空」。\src{B p.38} 尺寸的层级是全部推导的前提：$C$ 取得很小，$\Sigma$ 取得很大（$\Sigma$ 甚至可以是不紧的，例如 $\RR^2$）。因为扭化后的理论是拓扑场论，对度量没有依赖，所以把 $M$ 取成乘积、把两个因子的尺度拉开，都不损失一般性。

为什么这个层级正好给出二维理论？原因在 Kaluza--Klein 谱。把四维场按 $C$ 上的本征模式展开，第 $\mathbf n$ 个模式的质量平方约为 $\lambda_{\mathbf n}/R_C^2$；当能量远小于 $1/R_C$ 时，所有非零模式都被积掉，只剩零模。另一方面，零模沿 $\Sigma$ 可以缓变：场的 $\Sigma$ 方向导数每多一阶，作用量就多一个 $R_C^2$ 的因子，$\Sigma$ 很大时这一项可以任意小。于是「$C$ 小」保证只有零模参与低能物理，「$\Sigma$ 大」保证零模在 $\Sigma$ 上形成真正的二维场。两个条件缺一不可：若 $\Sigma$ 也小，就只剩量子力学；若 $C$ 不小，非零模式掺进来，二维描述失效。\src{B p.39}

还有一个超对称的账要算：四维的 16 个超荷在紧化后保留 8 个，它们分成左右各 4 个，即在 $\Sigma$ 上形成二维 $(4,4)$ 超对称。\src{B p.46} 这 8 个超荷正是后面二维 A-模型与 B-模型的超对称来源；若保留得更少，二维理论就不会是拓扑的。

\subsection{靶空间为什么是 $\MH$}

下一步要认出二维理论的靶空间。这要分两步，原文把它们压在一句话里，我们拆开看。

\textbf{第一步：认出超对称位形。}由第 3 讲，$t=\pm\ii$ 时四维的超对称位形是平坦 $G_{\mathbb C}$ 联络；更一般地，$t\neq0,\infty$ 时任何超对称位形都与 $t$ 无关，并且自动是平坦 $G_{\mathbb C}$ 联络（原文的定理 2）。\src{B pp.28--30} 所以在 $M=\Sigma\times C$ 上要研究平坦 $G_{\mathbb C}$ 丛 $E$。平坦丛在乘积流形上可以按两个因子分解：
\begin{equation}
E=\bigoplus_i E'_i\boxtimes E''_i,
\label{eq:E-decomp}
\end{equation}
其中 $E'_i$ 是 $\Sigma$ 上的平坦丛，$E''_i$ 是 $C$ 上的平坦丛。固定一点 $x\in\Sigma$，把 $E$ 限制到 $C_x=\{x\}\times C$ 上，得到
\begin{equation}
E_x=\bigoplus_i E'_i|_x\otimes E''_i,
\label{eq:E-x}
\end{equation}
它是 $C$ 上的一个平坦 $G_{\mathbb C}$ 丛；其同构类正是 $\MH(G,C)$ 的一个点（Hitchin 模空间与平坦丛模空间的对应，见第 1 讲）。让 $x$ 在 $\Sigma$ 上变动，就得到映射
\begin{equation}
\Phi:\ \Sigma\longrightarrow \MH(G,C),\qquad x\longmapsto [E_x].
\label{eq:Phi}
\end{equation}
这就是二维理论的场。原文在这里有一个技术约定：若某个 $E_x$ 是可约的平坦丛，它对应 $\MH$ 的奇异点；避开奇异点，即要求 $E$ 的限制不可约。\src{B p.30}

\textbf{第二步：认出能量。}光有映射还不够，还要说明低能作用量就是 sigma 模型的作用量。第 3 讲 §3.4 把四维玻色作用量写成了平方和加拓扑项（\src{B p.34 eq (3.47)--(3.48)}）：
\begin{equation}
S=\int_M\bigl(\text{平方和}\bigr)+\frac{\ii\Psi}{4\pi}\int_M\Tr(F\wedge F),
\label{eq:action-sum-sq}
\end{equation}
平方和为零正好是 Hitchin 方程
\begin{equation}
F-\varphi\wedge\varphi=0,\qquad D\varphi=0,\qquad D{*}\varphi=0,
\label{eq:hitchin}
\end{equation}
即 $A+\ii\varphi$ 平坦且 $D{*}\varphi=0$。在 $\Sigma$ 方向缓变的位形上，平方和项展开后正是映射 $\Phi$ 的动能（其二次项给出 $\MH$ 上的度量），拓扑项则给出二维的 $B$-场与 theta 角。于是低能有效理论就是以 $\MH(G,C)$ 为靶空间的二维 sigma 模型。原文 p.38 说这一步「由 §3.4 的结果得到」；把「认出靶空间」与「认出能量」分开，读者才不会误以为靶空间是从作用量直接读出来的。\src{B p.34}\src{B p.38}

顺带记下靶空间的几何：$\MH(G,C)$ 是超凯勒流形，复维数 $2(g-1)\dim G$，带一整族复结构（twistor 球面 $\PP^1_h$）。这个超凯勒结构是下一小节 A-模型/B-模型分类的全部来源。

\subsection{二维理论：A-模型、B-模型与 $\Psi$}

二维理论有 $(4,4)$ 超对称，靶空间 $\MH$ 又是超凯勒的，因此超对称变分可以按左右两个「扭化」独立选择复结构。设左右各选一个复结构 $J_+,J_-$（都来自 $\PP^1_h$），BPS 方程是
\begin{equation}
(1-\ii J_+)\,\bar\partial\Phi=0,\qquad
(1-\ii J_-)\,\partial\Phi=0.
\label{eq:bps2d}
\end{equation}
\src{B pp.59--61} 三种情形由此分出：

\begin{itemize}
\item $J_+=J_-=J$：两个方程合起来只依赖复结构 $J$，这是 \term{B-model}{B-模型}，其关联函数只依赖靶空间的复结构。
\item $J_+=-J_-=K$：方程只依赖辛形式，这是 \term{A-model}{A-模型}，其关联函数只依赖靶空间的辛结构（复化 Kähler 类）。
\item $J_-\neq\pm J_+$：既不是 A 也不是 B；但广义复几何告诉我们，它等价于带 $B$-场的 A-模型。原文只给结论、不给证明。\src{B p.59}\src{B pp.66--68}
\end{itemize}

四维的 $t$ 如何映到这一对复结构？原文给出
\begin{equation}
(w_+,w_-)=(-t,\ t^{-1}),
\label{eq:w-t}
\end{equation}
即左右复结构分别由 twistor 球面上的点 $w_+,w_-$ 指定；真正进入二维理论的是比值
\begin{equation}
q=\frac{w_+}{w_-}=-t^2.
\label{eq:q}
\end{equation}
\src{B p.61}\src{B p.62} 于是：$t=\pm\ii$ 时 $w_+=w_-$，二维理论是 \emph{B-模型}（$t=\ii$ 给复结构 $J$，$t=-\ii$ 给 $-J$）；$t$ 为实数时 $w_-=-1/\bar w_+$，二维理论是 \emph{A-模型}（$t=\pm1$ 给辛结构 $K$；$t=0,\infty$ 是退化极限）。这正是第 3 讲所说「$t=\pm\ii$ 像 B-模型、实 $t$ 像 A-模型」的精确版本。\src{B p.61}

但二维理论真正依赖的参数是 $\Psi$，不是 $t$。二维 A-模型的参数是复化 Kähler 类 $[B+\ii\omega]$；把四维的拓扑项降维后，可以算出
\begin{equation}
[B+\ii\omega]=-[\omega_I]\,\Psi,
\label{eq:B-i-omega}
\end{equation}
即在 $t\neq\pm\ii$ 时，$\PP^1$ 上每一点都等价于同一个 A-模型，只是 $B$-场平移了 $\Psi$ 倍。\src{B pp.68--69} 这就从二维一侧解释了「只有 $\Psi$ 有物理意义」。$t=\pm\ii$ 时上述推导因极点失效，二维理论是 B-模型，而 $\Psi=\infty$；此时理论对 $\tau$ 的无关性要单独论证，结论是：$t=\pm\ii$ 给出与 $\tau$ 无关的 B-模型。\src{B p.70}

把特殊值列成表（表~\ref{tab:psi-2d}），方向务必记清：
\begin{table}[htbp]
\centering
\caption{四维参数 $t$、典范参数 $\Psi$ 与二维理论基础模型的对应。$\Psi=\infty$ 对应 B-模型（复结构 $J$），$\Psi=0$ 对应 A-模型（辛结构 $K$），方向不要记反。}
\label{tab:psi-2d}
\begin{tabular}{@{}llll@{}}
\toprule
$t$ & $\Psi$ & 二维理论 & 备注 \\
\midrule
$\pm\ii$ & $\infty$ & B-模型，复结构 $J$（$t=-\ii$ 为 $-J$） & 与 $\tau$ 无关 \\
$\pm1$（且 $\operatorname{Re}\tau=0$） & $0$ & A-模型，辛结构对应 $K$ & 镜像对称的另一端 \\
其他实 $t$ & 实数 & A-模型 & 一般 A-模型 \\
$t=0,\infty$ & 退化 & 退化极限 & 非平凡瞬子在此附近 \\
一般 $t$ & 复数 & 广义复几何（等价于带 $B$-场的 A-模型） & 只有结论 \\
\bottomrule
\end{tabular}
\end{table}
\src{B pp.61--62}\src{B pp.68--71} 这里有一个方向容易记反，特别标出：\emph{$\Psi=\infty$ 对应 B-模型，$\Psi=0$ 对应 A-模型}；原文 p.70 正是从 $\Psi=\infty$ 的 B-模型出发做 $S$-对偶的。

\subsection{$S$-对偶 = 镜像对称：本讲核心}

现在到了本讲的核心。四维理论有 $S$-对偶：$G\leftrightarrow\Ldual$，$\tau\mapsto-1/\tau$。它在扭化族上的作用是
\begin{equation}
S:\ \Psi\longmapsto-\frac{1}{n_g\Psi},\qquad
T:\ \Psi\longmapsto\Psi+1,
\label{eq:S-T-2}
\end{equation}
矩阵形式为 $S=\begin{pmatrix}0&1/\sqrt{n_g}\\ -\sqrt{n_g}&0\end{pmatrix}$、$T=\begin{pmatrix}1&1\\0&1\end{pmatrix}$；$n_g$ 是第 3 讲介绍过的花边数，单边（ADE）时 $n_g=1$，$B,C,F_4$ 时 $n_g=2$，$G_2$ 时 $n_g=3$。\src{B p.70} 注意 $S$ 交换 $\Psi=0$ 与 $\Psi=\infty$：这正好把 A-模型端与 B-模型端互换。

把 $S$-对偶用在 $\Psi=\infty$ 的 B-模型上：对群 $G$，这是 $\MH(G)$ 上复结构 $J$ 的 B-模型；做 $S$-对偶后群变成 $\Ldual$，$\Psi$ 变成 $0$，二维理论变成 $\MH(\Ldual)$ 上的 A-模型。仔细跟踪复结构，可算出它对应的是辛结构 $K$，于是
\begin{equation}
\text{B-模型}_J\bigl(\MH(G)\bigr)\;\cong\;
\text{A-模型}_K\bigl(\MH(\Ldual)\bigr).
\label{eq:S-B-A}
\end{equation}
\src{B p.70} 这已经是一个镜像型陈述，但两边用的复结构不同。再做一次超凯勒旋转
\begin{equation}
\mathcal J:\ \varphi\longmapsto \ii\varphi,
\label{eq:J-rotation}
\end{equation}
它把 $K$ 转成 $J$，于是 \eqref{eq:S-B-A} 变成
\begin{equation}
\text{B-模型}_J\bigl(\MH(G)\bigr)\;\cong\;
\text{A-模型}_J\bigl(\MH(\Ldual)\bigr).
\label{eq:mirror-J}
\end{equation}
这正是「$\MH(G)$ 与 $\MH(\Ldual)$ 在复结构 $J$ 下互为镜像」的精确含义。$\mathcal J$ 不是扭化族的一员（它把 $\PP^1$ 上的点转出族外），原文说它「不十分自然」，但作为超凯勒流形上的旋转是合法的。\src{B pp.70--71}

为什么 $S$-对偶会给出镜像对称？几何上的原因是 Hitchin 纤维化。$\MH$ 有 Hitchin 纤维化 $\MH\to\mathcal B$，底 $\mathcal B$ 参数化 Higgs 场的特征多项式，纤维是阿贝尔簇（Jacobian 型环面）。$S$-对偶在底上作用平凡（单边时甚至 $\mathcal B\cong{}^L\mathcal B$），在纤维上则是 T-对偶：纤维与其对偶纤维互换。底加纤维的 T-对偶正是 SYZ 意义下的镜像对称。原文用 B-模型的局部算符验证了这一点：局部算符只测量底上的信息，B-模型的手征环等于全纯函数环，而全纯函数在紧纤维上是常数、故来自底，$S$-对偶给出 $R(G)\cong R(\Ldual)$、$\mathcal B\cong{}^L\mathcal B$；纤维上的零膜在 $S$-对偶下变成支撑在中间维数纤维上的 $(B,A,A)$-膜，其自由度正是对偶环面。\src{B pp.71--77} 这一整套是物理推导：原文明确说「相信」每个二维拓扑场论由两个广义复结构之一决定、与另一个无关，并说完整证明尚不存在（\src{B pp.66--67}）。SYZ 图像里的映射 $\Xi:\mathcal B\to{}^L\mathcal B$ 通常也假定为平凡映射。\src{B p.77}

把上面的两步画成交换图，就是本讲的核心等式（图 \ref{fig:core}）：
\begin{equation}
\text{4d TFT}_\Psi(G)\ \text{在}\ \Sigma\times C\ \text{上}
\quad\xrightarrow{\ \text{紧化}\ }\
\text{2d }\sigma\text{-模型，靶空间 }\MH(G).
\label{eq:core-eq}
\end{equation}
在 $\Psi=\infty$ 处，右边就是复结构 $J$ 上的 B-模型。
\begin{figure}[htbp]
\centering
\begin{tikzcd}[column sep=4.8em, row sep=3.8em]
  \text{4d TFT}_{\Psi=\infty}(G)\ \text{on}\ \Sigma\times C
    \arrow[r, "\text{紧化}"]
    \arrow[d, "\text{$S$-对偶}"']
  & \text{2d }B\text{-模型 on }(\MH(G),J)
    \arrow[d, "\text{镜像}"] \\
  \text{4d TFT}_{\Psi=0}(\Ldual)\ \text{on}\ \Sigma\times C
    \arrow[r, "\text{紧化}"']
  & \text{2d }A\text{-模型 on }(\MH(\Ldual),J)
\end{tikzcd}
\caption{本讲核心交换图。上排是群 $G$、$\Psi=\infty$ 的 B-模型；下排是 $\Ldual$、$\Psi=0$ 的 A-模型。左竖边是四维 $S$-对偶（$\Psi\mapsto-1/(n_g\Psi)$），右竖边是二维镜像对称。两条路径——先紧化再镜像，或先 $S$-对偶再紧化——给出同一个二维理论；交换图的右边用 \eqref{eq:J-rotation} 的 $\mathcal J$ 把辛结构 $K$ 换成 $J$。}
\label{fig:core}
\end{figure}

为什么两条路径给出同一个二维理论？物理上，$S$-对偶是四维理论的等价，紧化只是一个操作，二者可交换；所以交换图在四维一侧自动成立。真正有内容的是右边那条竖边：它把四维的 $S$-对偶翻译成了二维的镜像对称。换句话说，不是先证明镜像对称再用它，而是把 $S$-对偶紧化，得到镜像对称。这就是 Kapustin--Witten 的路线：几何 Langlands 的对偶性被读成四维规范理论的 $S$-对偶，而后者在紧化后变成二维镜像对称。\src{B pp.70--71}

最后照例提醒：这里的每一步都是物理推导，不是数学证明；原文只做非分歧情形；2024 年前后 Gaitsgory 等人的代数几何工作才给出几何 Langlands 的严格证明。\src{B p.77}

\begin{keybox}
紧化把四维扭化理论变成以 $\MH$ 为靶空间的二维 sigma 模型：$\Psi=\infty$ 时它是复结构 $J$ 的 B-模型，$\Psi=0$ 时它是辛结构 $K$ 的 A-模型，而四维 $S$-对偶把 B-模型换成 A-模型，于是 $\MH(G)$ 与 $\MH(\Ldual)$ 在 $J$ 下互为镜像。
\end{keybox}

\begin{follow}
\begin{itemize}
\item \textbf{B §4, pp.38--45}：紧化到 $\Sigma\times C$ 与靶空间 $\MH$ 的论证（重点 pp.38--39；pp.40--45 是 Hitchin 模空间的基本性质）。
\item \textbf{B §5.1, pp.56--61}：二维扭化理论、$J_+,J_-$ 与 A-模型/B-模型的判据。
\item \textbf{B §5.2, pp.61--68}：$t$ 到 $(w_+,w_-)$ 的映射、$\Psi$ 与广义复几何（本讲只取结论的部分）。
\item \textbf{B §5.3, pp.68--71}：$\Psi=\infty$ 的特化、$S$-对偶与镜像对称（本讲核心）。
\end{itemize}
\end{follow}

\begin{checkbox}
\begin{itemize}
\item 为什么 $C$ 必须小、$\Sigma$ 必须大？若把两个尺寸互换，低能理论会变成什么？
\item 为什么平坦 $G_{\mathbb C}$ 丛在每一点 $x\in\Sigma$ 的限制给出 $\MH$ 的点？$\MH$ 的奇点对应什么？
\item 用 \eqref{eq:bps2d} 说明 $J_+=J_-$ 与 $J_+=-J_-$ 分别给出 B-模型与 A-模型。
\item 验证 $t=\pm\ii$ 对应 $\Psi=\infty$、$t=\pm1$ 且 $\operatorname{Re}\tau=0$ 对应 $\Psi=0$；说明为什么这个方向容易记反。
\item 写出 $S$ 对 $\Psi$ 的作用与 $n_g$ 的定义；单边时 $S$ 是什么？
\item 交换图四条边各是什么？两条路径为什么给出同一个二维理论？
\end{itemize}
\end{checkbox}
